\section{Leaderboard submissions}
\label{sec:leaderboard}

\textbf{The uploaded system is a re-ranker, but not the one reported above, and the gap is not in
our favour.} Both are stated so that the better number never stands in for the uploaded one.

\textbf{Why the two differ.} The Codabench test files withhold \feat{article_ids_clicked}. We
checked this by comparing columns, not by assumption: the training behaviours file minus the test
behaviours file leaves exactly the outcome columns (clicks and the next read and scroll values).
Our five click-popularity features count clicks before each impression, and on the test period
those clicks \emph{are} the withheld labels, so the features cannot be computed there at any cost.
We replaced them with the five exposure features, which only need the in-view lists the test files
do contain. We also dropped \feat{bm25} from the submission for cost rather than availability. It
scores every article for every distinct history, which projects to about 1.6 hours on
\mind{}-large, and it is worth only $+$0.0004 test AUC in the ablation.

\begin{table}[tb]
\centering\small
\begin{tabular}{@{}lcc@{}}
\toprule
\eb{} small, our test split & features & per-impression AUC \\
\midrule
reported system & 22 & \EBFinalTest \\
uploaded system (no click popularity, no \feat{bm25}) & 16 & 0.7351 \\
\midrule
uploaded $-$ reported, paired bootstrap & & $-$0.0556 [$-$0.0565, $-$0.0547] \\
\bottomrule
\end{tabular}
\caption{The gap between what we report and what we upload, measured on our own labelled test
split. The \mind{} upload uses 11 features and scores 0.6460 on our \mind{} test split.}
\label{tab:submitgap}
\end{table}

\textbf{Scale and validation.} The \eb{} test set has 13,536,710 impressions and 205,925,868
candidate pairs; \mind{}-large has 2,370,727. Neither fits in memory, so prediction streams in
slices. Scoring \eb{} took 3\,h\,06\,m at 44\,GB peak on a cluster node, after three earlier
implementations were abandoned once their measured rates projected 38, 90 and 135 hours; batching
the LightGBM call, 13.5M invocations down to 541, was the fix that mattered. Each rewrite was
checked for byte-identical output. Because \mind{} allows one submission per day, every file
passes an offline validator first: exact inner filename (\texttt{prediction.txt} for \mind{},
\texttt{predictions.txt} for \eb{}), nothing else in the zip, row count, row order, and each row a
permutation of $1..n$. Row order matters most, since ranks are positional and a reordered file
scores as a plausible bad model rather than failing. The validator was itself checked by injecting
five violations and confirming each one fails.

\begin{table}[tb]
\centering\small
\begin{tabular}{@{}lcccc@{}}
\toprule
leaderboard & AUC & MRR & nDCG@5 & nDCG@10 \\
\midrule
\mind{} (Codabench 13967) & \pend{pending} & \pend{pending} & \pend{pending} & \pend{pending} \\
\eb{} (Codabench 2469) & \pend{pending} & \pend{pending} & \pend{pending} & \pend{pending} \\
\bottomrule
\end{tabular}
\caption{Official leaderboard scores. \pend{To fill in after upload; screenshots are in
Appendix~\ref{app:screenshots}.}}
\label{tab:leaderboard}
\end{table}

\textbf{Reading the official \mind{} MRR.} \mind{}'s scorer averages the reciprocal rank over every
clicked article, while the common definition credits only the first. The two agree when an
impression has one click. On our \mind{} test split, where 28.8\% of impressions have more than
one click, the definitions give 0.3108 against 0.2693 for the same ranking, with AUC and nDCG
unchanged. Our harness reports both, and the ``all clicks'' column is the one to compare with the
leaderboard.
